\documentclass[11pt,a4paper]{article}
\usepackage[margin=1in]{geometry}
\usepackage{graphicx}
\usepackage{booktabs}
\usepackage{amsmath}
\usepackage{hyperref}
\usepackage{xcolor}
\usepackage{float}

% TODO marker: pieces awaiting the in-flight train-set evals (trs2 batch) / wandb pull.
\newcommand{\TODO}[1]{\par\noindent\colorbox{yellow!40}{\textbf{TODO:} #1}\par}
\newcommand{\stdv}[1]{{\tiny$\pm$#1}}

\title{HumanEval Training Dynamics Analysis\\
\large RankAlign vs.\ TC methods --- why Upper and Multi differ}
\author{Automated Analysis Report (in progress)}
\date{June 19, 2026}

\begin{document}
\maketitle

\begin{abstract}
This report mirrors the earlier RankAlign training-dynamics analysis
(\texttt{analysis/report.tex}, membership/IFEval) but for the \textbf{HumanEval} task,
across \textbf{gemma-4-31b} and \textbf{qwen-3.5-9b}, comparing the two HumanEval datasets
\textbf{v2.1correct-upper} and \textbf{v2.1correct-multi}. It targets three questions:
\textbf{(Q1)} why does RankAlign (s2) beat our TC methods (s4 self-TC, s7 neg-TC) on
\emph{multi}, while our methods win on \emph{upper} (gemma-4)?
\textbf{(Q2)} is typicality-correction (TC) training helping?
\textbf{(Q3)} are the negative results on the test tables an \emph{overfitting} effect, or
are they already present on the train split?
All generator metrics are typicality-corrected (tc) gen ROC AUC unless noted; each method is
shown in its natural eval mode (RankAlign/s3 $\to$ Neg base; s4 $\to$ PMI base; s7 $\to$
Neg base), matching the convention in \texttt{docs/rerun\_only\_tables.tex}.
\end{abstract}

\noindent\textbf{Status (2026-06-19):} the TEST-set comparison (all settings, both models,
both datasets) is complete. The train-set dynamics currently exist only for \textbf{gemma-4
s2/s4} (GROUP B). Train-set evals for \textbf{gemma s3/s7} and \textbf{all qwen settings} are
running on mll (38-job \texttt{trs2} batch); sections depending on them are marked \TODO{}.

\tableofcontents
\newpage

%% ============================================================
\section{Q1: The Upper-vs-Multi Effect (TEST set)}
%% ============================================================

Table~\ref{tab:he_q1} reports tc gen ROC AUC ($\times100$) on the held-out test split for
every setting, each in its natural eval mode. The key pattern your question rests on:

\begin{itemize}
\item \textbf{gemma-4 / upper}: our methods \emph{win} --- FLORA-Neg (s7) $=94.8$ and
FLORA-PMI (s4) $=90.7$ both beat RankAlign (s2) $=88.2$.
\item \textbf{gemma-4 / multi}: RankAlign \emph{wins} --- s2 $=94.2$ vs s7 $=92.0$, s4 $=90.1$.
\item \textbf{qwen}: RankAlign wins on \emph{both} datasets (multi especially, $96.5$). The
``ours-wins-on-upper'' effect is \textbf{gemma-4-specific}.
\end{itemize}

\begin{table}[H]
\centering
\caption{HumanEval TEST tc gen ROC AUC ($\times100$). Each method in its natural eval mode
(Base: self; SFT/s4/s13: PMI base; RankAlign/New+fsx/s7: Neg base). gemma-4 has no base eval.}
\label{tab:he_q1}
\begin{tabular}{llccccccc}
\toprule
Model & Data & Base & SFT & RankAlign & New+fsx & FLORA-PMI (self-TC) & FLORA-Neg (neg-TC) & Consistency FT \\
\midrule
gemma-4-31b & upper & --- & 82.9 & 88.2 & 87.9 & 90.7 & 94.8 & 83.4 \\
gemma-4-31b & multi & --- & 87.9 & 94.2 & 94.0 & 90.1 & 92.0 & 87.3 \\
\midrule
qwen-3.5-9b & upper & 69.1 & 87.9 & 90.2 & 92.6 & 83.2 & 88.9 & 86.0 \\
qwen-3.5-9b & multi & 69.0 & 89.4 & 96.5 & 96.7 & 89.8 & 87.9 & 88.8 \\
\bottomrule
\end{tabular}

\end{table}

\TODO{Add the train-set counterpart of Table~\ref{tab:he_q1} once the \texttt{trs2} evals
land (gemma s3/s7 + all qwen). Then we can state, per setting, whether the upper/multi
ordering is the same on train as on test --- the crux of whether the effect is optimization
(seen on train) or generalization (only on test).}

%% ============================================================
\section{Q3: Train vs.\ Test --- Is It Overfitting?}
%% ============================================================

This is the central question. For the settings we already have on both splits
(\textbf{gemma-4 s2 RankAlign} and \textbf{s4 FLORA-PMI}, epoch 2), Table~\ref{tab:he_tt}
compares train-split and test-split tc gen ROC.

\begin{table}[H]
\centering
\caption{gemma-4 train vs.\ test tc gen ROC ($\times100$), epoch 2, each method in its
natural eval mode. Train = \texttt{--train} split, N=50 stratified candidates/problem.
$\Delta = \text{train}-\text{test}$ (large positive $\Rightarrow$ overfitting).}
\label{tab:he_tt}
\begin{tabular}{llccc}
\toprule
Data & Setting & Train & Test & $\Delta$(train$-$test) \\
\midrule
upper & RankAlign & 83.4 & 88.2 & -4.8 \\
upper & FLORA-PMI (self-TC) & 90.9 & 90.7 & 0.2 \\
multi & RankAlign & 94.6 & 94.2 & 0.4 \\
multi & FLORA-PMI (self-TC) & 99.0 & 90.1 & 8.9 \\
\bottomrule
\end{tabular}

\end{table}

\textbf{Early finding (gemma-4, s2 vs s4).} On \textbf{multi}, FLORA-PMI (s4) shows a large
train$\to$test drop ($99.0\to90.1$, $\Delta=+8.9$) while RankAlign generalizes almost
perfectly ($94.6\to94.2$, $\Delta=+0.4$). On \textbf{upper}, s4 does \emph{not} overfit
($90.9\to90.7$, $\Delta=+0.2$). This is consistent with the hypothesis that \textbf{our
TC method overfits on multi} --- it is best on the training candidates but loses the
advantage on held-out problems, which is exactly where RankAlign overtakes it at test time.
RankAlign's small train$\to$test gap suggests it learns a flatter, more transferable
ranking on this task.

\TODO{This only covers s2 vs s4. The method that \emph{wins} upper is \textbf{s7 (neg-TC)},
whose train-set eval is in the \texttt{trs2} batch --- add s7 (and s3, the no-TC baseline)
train-vs-test once done, for both gemma and qwen. Expected check: does s7 overfit on multi
(explaining its multi loss) but not on upper?}

\TODO{Add the qwen train-vs-test rows (s2/s3/s4/s7) once \texttt{trs2} lands --- needed to
confirm whether qwen's ``RankAlign wins both'' is also an overfitting story for the TC
methods, or something else.}

%% ============================================================
\section{Q2: Is TC Training Helping? (Focused TC comparison)}
%% ============================================================

The clean test of TC is s3 (New+fsx, \emph{no} TC) vs s4 (adds self-TC) and s7 (adds neg-TC),
evaluated in the matching TC mode (tc gen ROC, base-typ). Table~\ref{tab:he_tc} gives the
deltas at test (epoch 2).

\begin{table}[H]
\centering
\caption{TC effect at TEST: s3$\to$s4 (self mode) and s3$\to$s7 (neg mode), tc gen ROC
($\times100$); parenthetical = $\Delta$ from adding TC.}
\label{tab:he_tc}
\begin{tabular}{llcccc}
\toprule
Model & Data & s3 (self) & s4 (+self-TC) & s3 (neg) & s7 (+neg-TC) \\
\midrule
gemma-4-31b & upper & 87.5 & 90.7 (+3.2) & 87.9 & 94.8 (+6.8) \\
gemma-4-31b & multi & 90.9 & 90.1 (-0.7) & 94.0 & 92.0 (-2.1) \\
\midrule
qwen-3.5-9b & upper & 86.8 & 83.2 (-3.6) & 92.6 & 88.9 (-3.7) \\
qwen-3.5-9b & multi & 92.3 & 89.8 (-2.4) & 96.7 & 87.9 (-8.8) \\
\bottomrule
\end{tabular}

\end{table}

\textbf{Finding.} TC training \emph{helps only on gemma-4 / upper} (s4 $+3.2$, s7 $+6.8$ over
the no-TC s3). Everywhere else it \emph{hurts} at test: gemma multi ($-0.7$, $-2.1$), and qwen
on both datasets (down to $-8.8$ on qwen multi). This is the \emph{opposite} of the old
membership/IFEval report, where TC helped consistently --- on HumanEval, the TC signal
transfers only in the one regime (gemma upper) where our methods also win Table~\ref{tab:he_q1}.

\TODO{Add the train-side TC deltas once \texttt{trs2} lands. Hypothesis: TC's train-set gain
is large everywhere, but only transfers on gemma upper --- i.e.\ the test losses are the
overfitting signature of Section~2.}

%% ============================================================
\section{Generator--Validator Correlation (TEST)}
%% ============================================================

Spearman (Table~\ref{tab:he_sp}) and Pearson (Table~\ref{tab:he_pe}) between generator and
validator scores, epoch 2, each method in its natural eval mode (tc), mean $\pm$ SE over 82
problems.

\begin{table}[H]
\centering
\caption{Spearman $\rho$ (gen vs val), tc, $\times100$.}
\label{tab:he_sp}
\begin{tabular}{llccccccc}
\toprule
Model & Data & Base & SFT & RankAlign & New+fsx & FLORA-PMI (self-TC) & FLORA-Neg (neg-TC) & Consistency FT \\
\midrule
gemma-4-31b & upper & --- & 42.2\stdv{2.9} & 54.7\stdv{2.7} & 55.4\stdv{2.4} & 66.6\stdv{1.8} & 74.9\stdv{1.2} & 48.1\stdv{2.8} \\
gemma-4-31b & multi & --- & 62.2\stdv{1.7} & 75.9\stdv{1.2} & 68.9\stdv{1.7} & 73.0\stdv{1.4} & 79.3\stdv{1.2} & 58.8\stdv{1.7} \\
\midrule
qwen-3.5-9b & upper & 30.6\stdv{2.5} & 50.5\stdv{2.2} & 60.9\stdv{2.0} & 64.0\stdv{1.8} & 64.1\stdv{1.7} & 66.8\stdv{1.6} & 46.6\stdv{2.1} \\
qwen-3.5-9b & multi & 33.8\stdv{2.2} & 63.6\stdv{1.6} & 65.4\stdv{1.5} & 71.6\stdv{1.4} & 67.4\stdv{1.6} & 66.8\stdv{1.6} & 57.4\stdv{1.8} \\
\bottomrule
\end{tabular}

\end{table}

\begin{table}[H]
\centering
\caption{Pearson $r$ (gen vs val), tc, $\times100$.}
\label{tab:he_pe}
\begin{tabular}{llccccccc}
\toprule
Model & Data & Base & SFT & RankAlign & New+fsx & FLORA-PMI (self-TC) & FLORA-Neg (neg-TC) & Consistency FT \\
\midrule
gemma-4-31b & upper & --- & 31.8\stdv{3.2} & 55.5\stdv{2.8} & 60.2\stdv{2.5} & 63.3\stdv{1.8} & 75.8\stdv{1.3} & 37.8\stdv{3.1} \\
gemma-4-31b & multi & --- & 55.7\stdv{2.3} & 70.7\stdv{1.2} & 70.4\stdv{1.6} & 69.0\stdv{1.6} & 78.1\stdv{1.1} & 51.0\stdv{2.3} \\
\midrule
qwen-3.5-9b & upper & 26.5\stdv{2.4} & 45.4\stdv{2.3} & 69.5\stdv{1.7} & 67.6\stdv{1.8} & 58.9\stdv{2.0} & 68.5\stdv{1.5} & 41.7\stdv{2.3} \\
qwen-3.5-9b & multi & 27.1\stdv{2.4} & 57.8\stdv{1.9} & 70.1\stdv{1.1} & 72.4\stdv{1.4} & 66.0\stdv{1.7} & 67.9\stdv{1.4} & 53.1\stdv{2.0} \\
\bottomrule
\end{tabular}

\end{table}

\TODO{Generator--validator \textbf{concordance} (ordering-agreement fraction) and the
\textbf{scatter} / \textbf{score-delta histograms} need the raw per-candidate scores from mll
(\texttt{outputs\_gemma4\_mll\_tmp[-multi]/}, \texttt{outputs-rerun-wandb/}) --- pulling next.}

%% ============================================================
\section{Q2 (cont.): Training Diagnostics --- WandB}
%% ============================================================

WandB logging was \emph{online} for all HumanEval runs (project \texttt{rankalign}; run names
\texttt{g4-\{cu,cm\}-s\{S\}} and \texttt{q35-\{cu,cm\}-s\{S\}}).

\TODO{Pull and plot the loss curves (SFT vs RankAlign vs Ours; TC comparison) for gemma-4 and
qwen on upper and multi, mirroring report.tex Section 2. Also the loss-component breakdown
(Pref / NLL-G / NLL-V) --- check whether RankAlign's NLL-G stays bounded on HumanEval (it
exploded on IFEval in the old report). Reuse \texttt{analysis/scripts/analyze\_wandb\_training.py}.}

%% ============================================================
\section{Diagnostics: Concordance / Correlation / Scatter}
%% ============================================================

\TODO{Compute generator--validator concordance, Spearman, and Pearson (epoch 2) and the
gen-vs-val scatter / score-delta histograms from the raw candidate-level score CSVs on mll
(\texttt{outputs\_gemma4\_mll\_tmp[-multi]/}, \texttt{outputs-rerun-wandb/}, and the train
scores in \texttt{outputs-he-trainset/}). Mirror report.tex Sections 2.5--2.9. Note:
HumanEval has no category structure, so the per-category analysis is omitted.}

%% ============================================================
\section{Conclusions (preliminary)}
%% ============================================================

\begin{enumerate}
\item \textbf{The upper-vs-multi flip is real and gemma-4-specific}: our TC methods (s7, s4)
beat RankAlign on gemma-4 \emph{upper} but lose on \emph{multi}; on qwen RankAlign wins both.
\item \textbf{Preliminary overfitting signal}: gemma-4 FLORA-PMI overfits on multi
($\Delta=+8.9$ train$-$test) but not on upper ($\Delta=+0.2$), while RankAlign generalizes on
both. This is consistent with the multi loss being an overfitting effect.
\end{enumerate}

\TODO{Finalize conclusions once s7/s3/qwen train-set evals + WandB + concordance land. Open
questions: (a) does s7 overfit on multi too? (b) does TC's train-set advantage simply fail to
transfer on multi? (c) is qwen's behavior the same overfitting story or different?}

\paragraph{Reproducibility.} Tables built by
\texttt{analysis/scripts/build\_he\_report\_tables.py} from
\texttt{docs/he\_harvest\_2026-06-17/humaneval\_TEST\_metrics\_aggregated.csv} (test) and
\texttt{docs/he\_trainset\_2026-06-19/humaneval\_TRAIN\_metrics\_aggregated.csv} (train,
gemma s2/s4). trs2 train-set batch: \texttt{scripts/mll\_he\_trainset\_eval\_v2.sbatch} +
\texttt{launch\_he\_trainset\_extra.sh}; model origins in
\texttt{docs/he\_trainset\_eval\_model\_origins.md}.

\end{document}
